% TOP_PROBLEM: {"problemId": "problem.brauer-manin-obstruction-for-curves", "canonicalTitle": "Brauer-Manin obstruction for number-field curves: Hasse-principle form", "releaseRank": 210, "primaryDomain": "number_theory_arithmetic_geometry", "primaryDomainLabel": "Number theory & arithmetic geometry", "displayStatus": "open_with_solved_subcases", "releaseStatus": "open_with_solved_subcases"}
% !TeX program = lualatex
% Definition and sources only; this file is not an LLM solution attempt.
\documentclass[11pt]{article}
\usepackage[margin=25mm]{geometry}
\usepackage{fontspec}
\setmainfont{Latin Modern Roman}
\usepackage{amsmath,amssymb,mathtools}
\usepackage{unicode-math}
\setmathfont{Latin Modern Math}
\usepackage{xurl,hyperref}
\hypersetup{colorlinks=true,urlcolor=blue}
\setcounter{secnumdepth}{0}
\setlength{\parindent}{0pt}
\setlength{\parskip}{5pt}
\setlength{\emergencystretch}{3em}
\begin{document}
\section*{\texorpdfstring{Brauer-Manin obstruction for number-field curves: Hasse-principle form}{Brauer-Manin obstruction for number-field curves: Hasse-principle form}}\label{210}
\subsection{Definitions and mathematical statement}
Let $C/k$ be a smooth projective geometrically connected curve of genus at least two over a number field. Its Brauer--Manin set is
\[
 C({\mathbb{A}}_k)^{\operatorname{Br}}=
 \{(P_v):\sum_v\operatorname{inv}_v(\alpha(P_v))=0
 \text{ for every }\alpha\in H^2_{\mathrm{et}}(C,\mathbf G_m)\}.
\]
The local invariants take values in ${\mathbb{Q}}/{\mathbb{Z}}$ and the sum has only finitely many nonzero terms for fixed $\alpha,(P_v)$. The selected assertion is nonemptiness of this set implies $C(k)\ne\varnothing$. This is only the Hasse-principle form; density of rational points in the adelic set is not added. No rational divisor class of degree one or finiteness assumption on a Tate--Shafarevich group is permitted as an unstated premise.
\par\smallskip\textit{Formulation and concept references:} \href{https://arxiv.org/pdf/math/0606465v3}{[S1]}.

\subsection{Short English statement}
For a number-field curve of genus at least two, do adelic points passing every Brauer–Manin test guarantee an actual rational point?
\subsection{Sources}
\begin{itemize}
\item[S1]Michael Stoll, Finite descent obstructions and rational points on curves, arXiv math/0606465v3.
\url{https://arxiv.org/pdf/math/0606465v3}.
\textit{Formulation source.}
\end{itemize}
\end{document}
